% BHU PYQ -- Manually transcribed & error-corrected
\documentclass[11pt,a4paper]{article}

\usepackage[a4paper,top=2.5cm,bottom=2.5cm,left=2.8cm,right=2.8cm,headheight=14pt]{geometry}
\usepackage{amsmath,amssymb}
\usepackage[T1]{fontenc}
\usepackage[utf8]{inputenc}
\usepackage{lmodern}
\usepackage{microtype}
\usepackage{fancyhdr}
\usepackage{enumitem}
\usepackage{hyperref}
\usepackage{lastpage}
\usepackage{parskip}

\hypersetup{colorlinks=true,urlcolor=black,linkcolor=black,
  pdftitle={BSC-07A: Ancillary Physics-II -- BHU PYQ},pdfauthor={Banaras Hindu University}}

\pagestyle{fancy}\fancyhf{}
\renewcommand{\headrulewidth}{0.4pt}\renewcommand{\footrulewidth}{0.4pt}
\fancyhead[L]{\small   \textit{Banaras Hindu University}}
\fancyhead[R]{\small   \textit{BSC-07A: Ancillary Physics-II}}
\fancyfoot[C]{\small Page  \thepage\ of \pageref{LastPage}}


\newlist{parts}{enumerate}{2}
\setlist[parts,1]{label=(\alph*),leftmargin=*,itemsep=5pt,topsep=3pt}
\setlist[parts,2]{label=(\roman*),leftmargin=*,itemsep=3pt,topsep=2pt}

\newcommand{\pts}[1]{\hfill   \textit{[#1]}}


\begin{document}


\begin{center}
  {\large\bfseries BANARAS HINDU UNIVERSITY}\\[2pt]
  {
\normalsize B.Sc. (Hons.) Semester IV Examination, 2022-23}\\[6pt]
  \rule{0.8\linewidth}{0.5pt}\\[6pt]
  {\large\bfseries Paper No. BSC-07A: Ancillary Physics-II}\\[4pt]
  {
\normalsize Time: 3 Hours\hfill Maximum Marks: 70}
\end{center}
\rule{\linewidth}{0.4pt}
\smallskip



\noindent   \textit{Instructions: Answer   \textbf{five} questions including Question~1 which is compulsory. Symbols have their usual meanings.}
\bigskip




\noindent   \textbf{Question 1.} Answer any   \textbf{four} of the following: \pts{4 $ \times$ 3.5 = 14}

\begin{parts}
  \item State Fermat's principle of least time with an example.
  \item What is polarized light? What nature of light is shown by polarization?
  \item What will be the refractive index of a medium in which the speed of light is $2.5  \times 10^8\, \text{m/s}$?
  \item Give an innovative example to explain entropy.
  \item An ideal heat engine operating between the temperatures $T_1$ and $T_2$ ($T_1 > T_2$) has efficiency $\eta$. What would be the efficiency of the engine if you double the operating temperatures?
  \item Deduce the unit of Boltzmann's constant from the ideal gas equation.
\end{parts}

\medskip\rule{\linewidth}{0.2pt}




\noindent   \textbf{Question 2.}

\begin{parts}
  \item How do you define heat? Explain the heat capacity at constant volume and constant pressure. \pts{6}
  \item Write the formulae and use them to convert the boiling temperature of gold, $2966^\circ \text{C}$, into degrees Fahrenheit and Kelvin. \pts{6}
  \item What is the unit of pressure in the SI system of units? \pts{2}
\end{parts}

\medskip\rule{\linewidth}{0.2pt}




\noindent   \textbf{Question 3.}

\begin{parts}
  \item Define Viscosity and Thermal conductivity and give an example for each. \pts{6}
  \item The heat capacity of $0.125\, \text{kg}$ of water is measured to be $523\, \text{J/K}$ at room temperature. Calculate the heat capacity of water per unit mass and per unit volume. \pts{5}
  \item Which law of thermodynamics makes sense in the biological process of Photosynthesis? Explain. \pts{3}
\end{parts}

\medskip\rule{\linewidth}{0.2pt}




\noindent   \textbf{Question 4.}

\begin{parts}
  \item What is thermal equilibrium? Use it to explain the Zeroth law of thermodynamics. \pts{6}
  \item What is the relation between free energy and work? Starting from $F = U - TS$, where $F$ is the Helmholtz function and other notations carry the usual meaning, express the entropy and pressure at constant volume and temperature, respectively. \pts{6}
  \item What are the state variables and how do you label different thermodynamic processes by fixing one of these variables? \pts{2}
\end{parts}

\medskip\rule{\linewidth}{0.2pt}




\noindent   \textbf{Question 5.}

\begin{parts}
  \item Explain the differences between the following:
    
\begin{parts}
      \item Concave and Convex mirrors
      \item Interference and Diffraction
      \item Real image and Virtual image
    \end{parts} \pts{6}
  \item Explain Young's double-slit experiment with a neat ray diagram. Write the expression for the fringe-width of interference fringes. Why are all fringes of the same width? \pts{6}
  \item An object is kept at a position between the optical center and the focus of a convex lens. Show the image formation with the help of a neat diagram and write the properties of the image formed. \pts{2}
\end{parts}

\medskip\rule{\linewidth}{0.2pt}




\noindent   \textbf{Question 6.}

\begin{parts}
  \item Explain the optical phenomena involved in the following statements:
    
\begin{parts}
      \item Meters like ammeters and voltmeters use a mirror to avoid parallax error.
      \item A diamond sparkles more than a glass piece of the same shape and size.
      \item A cat's eye or road stud helps drivers by improving night visibility.
    \end{parts} \pts{6}
  \item In the dispersion of light by a prism, which color is more dispersed? Explain quantitatively using the formula of minimum deviation angle. \pts{5}
  \item If the critical angle of a medium is $30^\circ$, find the polarizing angle for the medium. \pts{3}
\end{parts}

\medskip\rule{\linewidth}{0.2pt}




\noindent   \textbf{Question 7.}

\begin{parts}
  \item In the figure shown below, when light falls on the vertical surface $BC$, it refracts along the direction $BC$. Find the incidence angle $i$. Given $\sin^{-1}(0.75) = 48.6^\circ$. \pts{5}
  \item What is farsightedness or hyperopia? How can it be removed? Explain with the help of a neat diagram. \pts{5}
  \item Find the critical angle for a ray of light passing from glass ($\mu_g = 1.5$) to water ($\mu_w = 1.33$). \pts{4}
\end{parts}

\vfill
\rule{\linewidth}{0.4pt}

\begin{center}\small
  \href{https://bhu-exam.vercel.app/}{Click here to visit our website}\\[4pt]
  \href{https://me-aryan.vercel.app/}{Click here to visit our contributor}\\[4pt]
  \href{https://quashan-me.vercel.app/}{Click here to visit our contributor}
\end{center}

\end{document}
